\scanpage{17}
\mspage{15}
\[
 I'_t=\prod_{\ell\ne p}\bigl(I'/[I',I']\bigr)(\ell)\simeq\prod_{\ell\ne p}\mathbb Z_\ell(1). \tag*{(3.\unsure{15|16}.1)}
\]

\deleted{Ceci posé, \unsure{on|nous} constate aussitôt que les opérations de \(\pi\)} \deleted{Bien entendu, \(I'\) opère trivialement \inserted{sur \deleted{\unsure{\(I\)|\(I'\)}} \(I'_{\mathrm{ab}}\)}, donc c'est le groupe \(\pi/I'\) qui opère sur \(I'_{\unsure{\mathrm{ab}|\mathrm{tab}}}\) et à fortiori sur \(I'_t\). Ceci posé, on constate que \(I/I'\) opère trivialement sur (3.16.1), et que les opérations de \(\pi/I\simeq\pi_0\) sur \(I_t\) ainsi}

\deleted{\unsure{Or|Et}} l'isomorphisme \deleted{de \(\pi\)} (3.16.1) respecte les opérations de \(\pi\) sur les deux mbs, où l'opération sur le deuxième mb. \deleted{\unsure{est|se}} se fait via \(\pi_0\) comme expliqué déjà plus haut.

Ceci ne fait qu'exprimer le fait que l'isom canonique (2\,15.1, est compatible avec ``transport de structure''.

En particulier, on trouve que \(I\deleted{/\unsure{I'|I}}\) opère trivialement sur \(I'_t\), en d'autres termes

\uline{Corollaire} 3.1\deleted{\unsure{7|9}}6. Pour tout ss-groupe \inserted{ouvert} invariant \(I'\) \deleted{d'indice \unsure{fini|fin.}} de \(I\), l'extension naturelle de \(I/I'\) par le quotient \(I'_t\simeq\prod_{\ell\ne p}\mathbb Z_\ell(1)\) de \(I'\) est une extension \uline{centrale}.

% archive sheet number: 16
